% =========================================================
% PQ-BLE-HANDSHAKE - Tesina LaTeX pronta da compilare
% =========================================================
% !TEX program = latexmk
% Compilazione consigliata:
%   latexmk -pdf -interaction=nonstopmode tesina_pq_ble_handshake.tex

\documentclass[11pt,a4paper]{article}

% =========================================================
% METADATI DEL PROGETTO
% =========================================================
\newcommand{\Course}{Sicurezza nelle Reti}
\newcommand{\CourseAcronym}{PQ-BLE-HANDSHAKE}
\newcommand{\Professor}{Stelvio Cimato}
\newcommand{\University}{Universit\`a degli Studi di Milano - Sicurezza Informatica}
\newcommand{\AcademicYear}{2025/2026}
\newcommand{\Student}{Alessio Bonora}
\newcommand{\StudentEmail}{alebonora2003@gmail.com}
\newcommand{\ProjectTitle}{Canale Sicuro Post-Quantum su Bluetooth Low Energy}
\newcommand{\ProjectSubtitle}{Handshake GATT-based con ML-KEM-768 e SAS Numeric Comparison}
\newcommand{\ProjectDate}{10 luglio 2026}
\newcommand{\ProjectRepo}{\url{https://github.com/AleBonora3/pq-ble-handshake/tree/main/}}

% =========================================================
% LINGUA, CODIFICA, FONT
% =========================================================
\usepackage[utf8]{inputenc}
\usepackage[T1]{fontenc}
\usepackage[italian]{babel}
\usepackage{lmodern}
\usepackage[scaled=0.95]{inconsolata}
\usepackage{microtype}
\usepackage{textcomp}
\usepackage{eurosym}

% =========================================================
% LAYOUT E TESTO
% =========================================================
\usepackage[a4paper,margin=1.85cm]{geometry}
\usepackage{setspace}
\setstretch{1.08}
\setlength{\parindent}{0pt}
\setlength{\parskip}{3pt plus 1pt minus 1pt}
\usepackage{enumitem}
\setlist{topsep=2pt,partopsep=0pt,itemsep=1pt,parsep=0pt,leftmargin=1.6em}

% =========================================================
% COLORI
% =========================================================
\usepackage{xcolor}
\definecolor{Primary}{HTML}{1F4E79}
\definecolor{Secondary}{HTML}{2E7D32}
\definecolor{Accent}{HTML}{8E24AA}
\definecolor{Warning}{HTML}{C62828}
\definecolor{SoftBg}{HTML}{F6F8FA}
\definecolor{SoftBorder}{HTML}{D0D7DE}
\definecolor{CodeBg}{HTML}{F7F9FB}
\definecolor{CodeBorder}{HTML}{D9E2EC}
\definecolor{CodeKeyword}{HTML}{7B1FA2}
\definecolor{CodeString}{HTML}{0B8043}
\definecolor{CodeComment}{HTML}{546E7A}
\definecolor{CodeNumber}{HTML}{1565C0}
\definecolor{CodeIdentifier}{HTML}{212121}
\definecolor{InlineBg}{HTML}{EEF3F8}
\definecolor{CodeText}{HTML}{212121}

% =========================================================
% FIGURE E TABELLE
% =========================================================
\usepackage{graphicx}
\graphicspath{{images/}{../docs/images/}}

% Cerca prima nella cartella locale images/, poi nella cartella della repo
% ../docs/images/. Se il file manca, mostra un segnaposto compilabile.
\newcommand{\safeincludegraphics}[2][width=\linewidth]{%
  \IfFileExists{images/#2}{%
    \includegraphics[#1]{images/#2}%
  }{%
    \IfFileExists{../docs/images/#2}{%
      \includegraphics[#1]{../docs/images/#2}%
    }{%
      \fbox{\parbox[c][5.2cm][c]{0.92\linewidth}{%
        \centering\small\textcolor{black!60}{Immagine non trovata.\\[0.4em]
        File atteso: \texttt{\detokenize{#2}} in \texttt{images/} oppure \texttt{../docs/images/}.}%
      }}%
    }%
  }%
}
\usepackage{float}
\usepackage{caption}
\usepackage{subcaption}
\captionsetup{font=small,labelfont=bf,skip=4pt}
\usepackage{booktabs}
\usepackage{longtable}
\usepackage{array}
\usepackage{multirow}
\usepackage{makecell}
\usepackage{tabularx}
\usepackage{calc}

% Pandoc longtable helpers
\usepackage{iftex}
\usepackage{etoolbox}
\providecommand{\tightlist}{\setlength{\itemsep}{0pt}\setlength{\parskip}{0pt}}
\newcommand{\passthrough}[1]{#1}
\setlength{\emergencystretch}{3em}

% =========================================================
% MATEMATICA
% =========================================================
\usepackage{amsmath,amssymb,mathtools}

% =========================================================
% RISULTATI BENCHMARK IMPORTATI DA JSON
% =========================================================
% Il file benchmark_values.tex viene generato da
% scripts/generate_benchmark_latex.py a partire dagli artefatti JSON reali.
\newif\ifbenchmarkdata
\benchmarkdatafalse
\providecommand{\BenchGeneratedAt}{non disponibile}
\providecommand{\BenchIterations}{N/D}
\providecommand{\BenchKeygenMean}{N/D}
\providecommand{\BenchKeygenMedian}{N/D}
\providecommand{\BenchKeygenPnn}{N/D}
\providecommand{\BenchEncapsMean}{N/D}
\providecommand{\BenchEncapsMedian}{N/D}
\providecommand{\BenchEncapsPnn}{N/D}
\providecommand{\BenchDecapsMean}{N/D}
\providecommand{\BenchDecapsMedian}{N/D}
\providecommand{\BenchDecapsPnn}{N/D}
\providecommand{\BenchSasMean}{N/D}
\providecommand{\BenchHkdfMean}{N/D}
\providecommand{\BenchTotalMeanMs}{N/D}
\providecommand{\BenchTotalMedianMs}{N/D}
\providecommand{\BenchTotalPnnMs}{N/D}
\providecommand{\BenchHandshakeRate}{N/D}
\providecommand{\BenchEncSixtyFour}{N/D}
\providecommand{\BenchDecSixtyFour}{N/D}
\providecommand{\BenchEncTwoFiftySix}{N/D}
\providecommand{\BenchDecTwoFiftySix}{N/D}
\providecommand{\BenchEncFiveTwelve}{N/D}
\providecommand{\BenchDecFiveTwelve}{N/D}
\providecommand{\BenchEncOneK}{N/D}
\providecommand{\BenchDecOneK}{N/D}
\providecommand{\BenchEncFourK}{N/D}
\providecommand{\BenchDecFourK}{N/D}
\providecommand{\BenchEncSixteenK}{N/D}
\providecommand{\BenchDecSixteenK}{N/D}
\providecommand{\BenchFragMtuTwoFourSevenUs}{N/D}
\providecommand{\BenchReasmMtuTwoFourSevenUs}{N/D}
\providecommand{\BenchFragMtuFiveTwelveUs}{N/D}
\providecommand{\BenchReasmMtuFiveTwelveUs}{N/D}
\IfFileExists{benchmark_values.tex}{% =========================================================
% RISULTATI BENCHMARK REALI
% Generato automaticamente dagli artefatti JSON in benchmarks/results/
% Non modificare manualmente: eseguire scripts/generate_benchmark_latex.py
% =========================================================

\benchmarkdatatrue

\renewcommand{\BenchGeneratedAt}{21 luglio 2026, ore 14:20 UTC}

% Configurazioni dei diversi benchmark
\renewcommand{\BenchIterations}{1000}
\renewcommand{\BenchCryptoIterations}{1000}
\renewcommand{\BenchThroughputTrials}{5}
\renewcommand{\BenchFragIterations}{10000}

% ML-KEM-768 / SAS / HKDF: tempi in microsecondi
\renewcommand{\BenchKeygenMean}{161,63}
\renewcommand{\BenchKeygenMedian}{123,75}
\renewcommand{\BenchKeygenPnn}{562,30}

\renewcommand{\BenchEncapsMean}{161,03}
\renewcommand{\BenchEncapsMedian}{125,60}
\renewcommand{\BenchEncapsPnn}{527,10}

\renewcommand{\BenchDecapsMean}{74,60}
\renewcommand{\BenchDecapsMedian}{59,10}
\renewcommand{\BenchDecapsPnn}{215,20}

\renewcommand{\BenchSasMean}{8,54}
\renewcommand{\BenchSasMedian}{5,50}
\renewcommand{\BenchSasPnn}{34,40}

\renewcommand{\BenchHkdfMean}{14,55}
\renewcommand{\BenchHkdfMedian}{9,10}
\renewcommand{\BenchHkdfPnn}{68,60}

% Totale crittografico in millisecondi
\renewcommand{\BenchTotalMeanMs}{0,4204}
\renewcommand{\BenchTotalMedianMs}{0,3260}
\renewcommand{\BenchTotalPnnMs}{1,4054}
\renewcommand{\BenchHandshakeRate}{2\,378,9}

% Dimensioni del materiale crittografico ML-KEM-768
\renewcommand{\BenchPkBytes}{1184}
\renewcommand{\BenchSkBytes}{2400}
\renewcommand{\BenchCtBytes}{1088}
\renewcommand{\BenchSsBytes}{32}
\renewcommand{\BenchKemExchangeBytes}{2272}

% Throughput AES-256-GCM in KiB/s
\renewcommand{\BenchWireOverheadBytes}{37}

\renewcommand{\BenchEncSixtyFour}{31\,976}
\renewcommand{\BenchDecSixtyFour}{28\,634}
\renewcommand{\BenchWireSixtyFour}{101}
\renewcommand{\BenchOverheadSixtyFourPct}{57,81}

\renewcommand{\BenchEncTwoFiftySix}{125\,437}
\renewcommand{\BenchDecTwoFiftySix}{105\,418}
\renewcommand{\BenchWireTwoFiftySix}{293}
\renewcommand{\BenchOverheadTwoFiftySixPct}{14,45}

\renewcommand{\BenchEncFiveTwelve}{239\,123}
\renewcommand{\BenchDecFiveTwelve}{193\,878}
\renewcommand{\BenchWireFiveTwelve}{549}
\renewcommand{\BenchOverheadFiveTwelvePct}{7,23}

\renewcommand{\BenchEncOneK}{465\,468}
\renewcommand{\BenchDecOneK}{386\,306}
\renewcommand{\BenchWireOneK}{1061}
\renewcommand{\BenchOverheadOneKPct}{3,61}

\renewcommand{\BenchEncFourK}{1\,193\,458}
\renewcommand{\BenchDecFourK}{1\,057\,950}
\renewcommand{\BenchWireFourK}{4133}
\renewcommand{\BenchOverheadFourKPct}{0,90}

\renewcommand{\BenchEncSixteenK}{2\,628\,523}
\renewcommand{\BenchDecSixteenK}{1\,959\,725}
\renewcommand{\BenchWireSixteenK}{16421}
\renewcommand{\BenchOverheadSixteenKPct}{0,23}

% =========================================================
% FRAMMENTAZIONE APPLICATIVA DEL CIPHERTEXT ML-KEM-768
% =========================================================
\renewcommand{\BenchFragHeaderBytes}{4}

% Configurazione usata nel PoC validato
\renewcommand{\BenchFragTwoFourSevenLogicalSize}{247}
\renewcommand{\BenchFragTwoFourSevenPayload}{243}
\renewcommand{\BenchFragTwoFourSevenCount}{5}
\renewcommand{\BenchFragTwoFourSevenWireBytes}{1108}
\renewcommand{\BenchFragTwoFourSevenOverheadBytes}{20}
\renewcommand{\BenchFragTwoFourSevenOverheadPct}{1,84}
\renewcommand{\BenchFragMtuTwoFourSevenUs}{4,79}
\renewcommand{\BenchReasmMtuTwoFourSevenUs}{6,04}

% Configurazione di confronto
\renewcommand{\BenchFragFiveTwelveLogicalSize}{512}
\renewcommand{\BenchFragFiveTwelvePayload}{508}
\renewcommand{\BenchFragFiveTwelveCount}{3}
\renewcommand{\BenchFragFiveTwelveWireBytes}{1100}
\renewcommand{\BenchFragFiveTwelveOverheadBytes}{12}
\renewcommand{\BenchFragFiveTwelveOverheadPct}{1,10}
\renewcommand{\BenchFragMtuFiveTwelveUs}{3,46}
\renewcommand{\BenchReasmMtuFiveTwelveUs}{4,48}

% Public key trasportata tramite ATT Long Read / Read Blob
\renewcommand{\BenchPkTransport}{ATT Long Read / Read Blob}
\renewcommand{\BenchFullHandshakeAppBytes}{2292}

% =========================================================
% SESSION RESUMPTION LATO PYTHON
% =========================================================
\renewcommand{\BenchResumeRequestBytes}{21}
\renewcommand{\BenchResumeResponseBytes}{9}
\renewcommand{\BenchResumeExchangeBytes}{30}
}{}


% =========================================================
% SEZIONI E INDICE
% =========================================================
\usepackage{titlesec}
\usepackage{tocloft}
\titleformat{\section}{\Large\bfseries\color{Primary}}{\thesection.}{0.5em}{}
\titleformat{\subsection}{\large\bfseries\color{Primary}}{\thesubsection.}{0.5em}{}
\titleformat{\subsubsection}{\normalsize\bfseries\color{Primary}}{\thesubsubsection.}{0.5em}{}
\titleformat{\paragraph}[runin]{\normalsize\bfseries\color{Accent}}{\theparagraph.}{0.5em}{}[:]
\titlespacing*{\section}{0pt}{1.4ex plus .2ex}{0.7ex}
\titlespacing*{\subsection}{0pt}{1.1ex plus .2ex}{0.5ex}
\titlespacing*{\subsubsection}{0pt}{0.9ex plus .2ex}{0.4ex}
\titlespacing*{\paragraph}{0pt}{0.6ex plus .2ex}{0.5em}
\setcounter{secnumdepth}{3}
\setcounter{tocdepth}{2}

% =========================================================
% BOX / CALLOUT
% =========================================================
\usepackage[most]{tcolorbox}
\tcbset{enhanced,boxrule=0.6pt,arc=2pt,left=6pt,right=6pt,top=5pt,bottom=5pt,before skip=6pt,after skip=6pt}
\newtcolorbox{callout}{title=Nota,colback=SoftBg,colframe=Primary!55,coltitle=black,fonttitle=\bfseries,borderline west={2pt}{0pt}{Primary!70}}
\newtcolorbox{importantbox}{title=Importante,colback=Warning!3,colframe=Warning!50,coltitle=black,fonttitle=\bfseries,borderline west={2pt}{0pt}{Warning!70}}

% =========================================================
% CODICE (LISTINGS)
% =========================================================
\usepackage{listings}
\usepackage{upquote}
\newcommand{\code}[1]{{\setlength{\fboxsep}{2pt}\colorbox{InlineBg}{\texttt{#1}}}}
\lstset{inputencoding=utf8,extendedchars=true,upquote=true,showstringspaces=false,showtabs=false,showspaces=false,breaklines=true,breakatwhitespace=false,tabsize=2,keepspaces=true,columns=flexible,literate={à}{{\`a}}1 {è}{{\`e}}1 {é}{{\'e}}1 {ì}{{\`i}}1 {ò}{{\`o}}1 {ù}{{\`u}}1 {À}{{\`A}}1 {È}{{\`E}}1 {É}{{\'E}}1 {Ì}{{\`I}}1 {Ò}{{\`O}}1 {Ù}{{\`U}}1 {’}{{'}}1 {‘}{{'}}1 {“}{{``}}1 {”}{{''}}1}
\lstdefinestyle{codestyle}{backgroundcolor=\color{CodeBg},frame=single,rulecolor=\color{CodeBorder},frameround=tttt,framesep=5pt,xleftmargin=8pt,xrightmargin=4pt,aboveskip=7pt,belowskip=7pt,basicstyle=\ttfamily\small\color{CodeText},keywordstyle=\bfseries\color{CodeKeyword},stringstyle=\color{CodeString},commentstyle=\itshape\color{CodeComment},identifierstyle=\color{CodeIdentifier},numbers=none,captionpos=b}
\lstdefinestyle{pythonstyle}{style=codestyle,language=Python,morekeywords={self,True,False,None,match,case,async,await}}
\lstdefinestyle{cstyle}{style=codestyle,language=C}
\lstdefinestyle{bashstyle}{style=codestyle,language=bash,morekeywords={sudo,cd,ls,pwd,mkdir,rm,cp,mv,chmod,chown,grep,find,cat,echo,kill,killall,ps,top,gcc,g++,make,nano,vim,python,pip,west}}
\lstset{style=codestyle}

% =========================================================
% HEADER E FOOTER
% =========================================================
\usepackage{fancyhdr}
\pagestyle{fancy}
\setlength{\headheight}{14pt}
\fancyhf{}
\fancyhead[L]{\small\textcolor{Primary}{\CourseAcronym}}
\fancyhead[C]{\small\Course}
\fancyhead[R]{\small AA \AcademicYear}
\fancyfoot[L]{\small UNIMI - Sicurezza Informatica}
\fancyfoot[C]{\small\Student}
\fancyfoot[R]{\small\thepage}
\renewcommand{\headrulewidth}{0.4pt}
\renewcommand{\footrulewidth}{0.4pt}

% =========================================================
% LINK E PDF
% =========================================================
\usepackage{xurl}
\usepackage[hidelinks]{hyperref}
\usepackage{bookmark}

% =========================================================
% FRONTESPIZIO
% =========================================================
\newcommand{\makeprojecttitle}{%
\begin{titlepage}
  \centering
  \vspace*{1.5cm}
  {\Huge\bfseries\color{Primary}\ProjectTitle\par}
  \vspace{0.5cm}
  {\Large\bfseries\ProjectSubtitle\par}
  \vspace{1.4cm}
  {\large Tesina per il corso di \textbf{\Course} (6 CFU)\par}
  \vspace{0.2cm}
  {\large Laurea Magistrale - A.A. \AcademicYear\par}
  \vspace{1.2cm}
  \begin{tcolorbox}[width=0.82\textwidth,colback=SoftBg,colframe=Primary!45]
    \begin{tabularx}{\textwidth}{@{}lX@{}}
      \textbf{Autore:} & \Student \\
      \textbf{Email:} & \StudentEmail \\
      \textbf{Docente:} & \Professor \\
      \textbf{Universit\`a:} & \University \\
      \textbf{Data:} & \ProjectDate \\
      \textbf{Github:} & \ProjectRepo \\
    \end{tabularx}
  \end{tcolorbox}
  \vfill
  {\small Proof-of-concept di handshake post-quantum sopra BLE GATT con validazione su nRF54L15 DK.\par}
\end{titlepage}
}

\begin{document}
\makeprojecttitle
\pagenumbering{roman}
\section*{Abstract}
\addcontentsline{toc}{section}{Abstract}

Il Bluetooth Low Energy (BLE) protegge le comunicazioni mediante i
meccanismi previsti dal Security Manager, tra cui LE Secure Connections
basato su ECDH P-256. Tale soluzione è adeguata nel modello di minaccia
classico, ma non è progettata per resistere a un futuro calcolatore
quantistico crittograficamente rilevante. In uno scenario
\emph{store-now-decrypt-later}, un attaccante può registrare oggi il
traffico e tentare di decifrarlo in futuro.

Questa tesina presenta \textbf{PQ-BLE-HANDSHAKE}, un proof-of-concept di
handshake post-quantum a livello applicativo sopra BLE GATT, senza
modificare lo stack Bluetooth nativo. Il protocollo combina
\textbf{ML-KEM-768}, \textbf{SAS Numeric Comparison} a sei cifre,
\textbf{HKDF-SHA256} e un canale \textbf{AES-256-GCM} con Additional
Authenticated Data (AAD), sequence number monotono e protezione da replay.
La session resumption è limitata temporalmente a 24 ore e a 100 resume
riusciti; timeout e tentativi falliti non consumano il contatore.

La validazione è articolata su tre livelli. La logica crittografica e di
protocollo è coperta da una suite attiva di \textbf{101 test automatizzati}.
Un firmware Zephyr per \textbf{nRF54L15 DK} è stato compilato, flashato e
usato come peripheral reale. Il firmware espone una \textbf{public key
ML-KEM-768 valida, generata offline}; il central ne verifica il fingerprint
SHA-256, esegue l'incapsulamento e trasferisce il ciphertext da 1088 byte in
cinque frammenti con MTU 247. Infine, una cattura passiva con
\textbf{nRF52840 Dongle} e Wireshark documenta MTU exchange, long read della
public key, trasferimento del ciphertext, comando \texttt{START} e
notification finale da 57 byte.

Il limite principale è dichiarato esplicitamente: la secret key
corrispondente non è memorizzata sul DK e il firmware non esegue ancora
ML-KEM decapsulation, HKDF o AES-GCM on-chip. La prova hardware dimostra
quindi il trasporto e il riassemblaggio reali su BLE/GATT, mentre la
completezza crittografica è verificata lato Python. I benchmark sono
riproducibili e importati nella relazione dagli artefatti JSON prodotti
dalla suite dedicata.

\clearpage
\tableofcontents
\clearpage
\pagenumbering{arabic}

\section{Introduzione}\label{capitolo-1---introduzione}

\subsection{Continuità con il percorso di
ricerca}\label{continuituxe0-con-il-percorso-di-ricerca}

Questo progetto si inserisce in un percorso di ricerca sulla sicurezza
Bluetooth Low Energy iniziato nella mia tesi triennale, dedicata
all'analisi della sicurezza BLE tramite sniffing passivo con nRF52840
Dongle, Wireshark e una pipeline di analisi dei pacchetti.

La tesi precedente aveva un obiettivo prevalentemente diagnostico:
osservare il pairing BLE, analizzare il traffico SMP, verificare la
presenza o meno di protezioni MITM e valutare i limiti della sicurezza
BLE tradizionale. Questo progetto rappresenta il passaggio successivo:
invece di limitarsi a evidenziare il problema, propone e valida
sperimentalmente una possibile soluzione applicativa post-quantum.

\begin{lstlisting}
TESI TRIENNALE                         PROGETTO MAGISTRALE
Analisi sicurezza BLE                  PQ-BLE-HANDSHAKE

Sniffing passivo pairing        ->      Handshake post-quantum applicativo
Analisi ECDH/SMP                ->      ML-KEM-768 sopra GATT
nRF52840 come sniffer           ->      nRF52840 per validazione Wireshark
Pipeline di audit BLE           ->      Validazione trasporto BLE/GATT reale
\end{lstlisting}

Il progetto riutilizza quindi sia la strumentazione sia la conoscenza
acquisita: il nRF52840 Dongle viene usato per osservare il traffico BLE
e verificare che il nuovo protocollo venga trasportato come normali
operazioni ATT/GATT.

\subsection{Contesto}\label{contesto}

BLE è largamente usato in dispositivi IoT, wearable, sensori, smart home
e dispositivi medicali. La sicurezza BLE è quindi rilevante non solo per
la privacy, ma anche per scenari in cui i dati possono avere valore a
lungo termine.

Il pairing BLE moderno, LE Secure Connections, usa ECDH P-256 per
stabilire chiavi condivise. Questa scelta è sicura contro attaccanti
classici, ma non contro un quantum computer sufficientemente potente. La
conseguenza è che traffico registrato oggi potrebbe essere decrittato in
futuro.

\subsection{Minaccia
store-now-decrypt-later}\label{minaccia-store-now-decrypt-later}

Nello scenario \emph{store-now-decrypt-later}, l'attaccante non deve
rompere la crittografia immediatamente. È sufficiente registrare
traffico cifrato oggi e attendere che diventi disponibile una tecnologia
in grado di rompere ECDH. Questo modello è particolarmente rilevante per
dati sanitari, credenziali, segreti industriali e informazioni personali
con lunga vita utile.

\subsection{Stato dell'arte e
motivazione}\label{stato-dellarte-e-motivazione}

La migrazione alla crittografia post-quantum è già iniziata in
protocolli come TLS, Signal e sistemi di messaggistica moderni. BLE,
invece, non ha ancora una modalità post-quantum standardizzata nel
Security Manager. Per questo motivo il progetto adotta un approccio
pragmatico: non modifica il pairing BLE, ma costruisce un canale sicuro
post-quantum a livello applicativo sopra GATT.

L'analogia è con protocolli come HTTPS: il trasporto sottostante non
viene modificato, ma la sicurezza applicativa viene costruita sopra di
esso.

\subsection{Obiettivi del progetto}\label{obiettivi-del-progetto}

Gli obiettivi sono:

\begin{enumerate}
\def\labelenumi{\arabic{enumi}.}
\tightlist
\item
  progettare un handshake post-quantum compatibile con BLE GATT;
\item
  implementare un proof-of-concept modulare in Python;
\item
  integrare session key derivation, SAS, canale AES-GCM e session
  resumption;
\item
  preparare e validare un firmware Zephyr su nRF54L15 DK per il
  trasporto BLE/GATT;
\item
  catturare il traffico con nRF52840 Dongle e Wireshark;
\item
  documentare limiti, risultati e sviluppi futuri in modo chiaro.
\end{enumerate}

\begin{center}\rule{0.5\linewidth}{0.5pt}\end{center}

\section{Background}\label{capitolo-2---background}

\subsection{BLE, ATT e GATT}\label{ble-att-e-gatt}

Lo stack BLE può essere schematizzato così:

\begin{lstlisting}
APPLICATION      <- PQ-BLE-HANDSHAKE
GATT / ATT       <- trasporto usato dal progetto
SMP              <- non usato nel firmware demo
L2CAP
Link Layer / PHY
\end{lstlisting}

Il Generic Attribute Profile (GATT) organizza i dati in \textbf{Service}
e \textbf{Characteristic}. Ogni characteristic ha un UUID, proprietà
come READ/WRITE/NOTIFY e un valore trasportato tramite ATT. Nel progetto
PQ-BLE, GATT viene usato come trasporto standard per scambiare materiale
ML-KEM e messaggi di controllo.

\subsection{Ruoli BLE}\label{ruoli-ble}

Nel progetto:

\begin{longtable}[]{@{}
  >{\raggedright\arraybackslash}p{(\columnwidth - 4\tabcolsep) * \real{0.3333}}
  >{\raggedright\arraybackslash}p{(\columnwidth - 4\tabcolsep) * \real{0.3333}}
  >{\raggedright\arraybackslash}p{(\columnwidth - 4\tabcolsep) * \real{0.3333}}@{}}
\toprule\noalign{}
\begin{minipage}[b]{\linewidth}\raggedright
Ruolo
\end{minipage} & \begin{minipage}[b]{\linewidth}\raggedright
Descrizione
\end{minipage} & \begin{minipage}[b]{\linewidth}\raggedright
Implementazione
\end{minipage} \\
\midrule\noalign{}
\endhead
\bottomrule\noalign{}
\endlastfoot
Central / GATT Client & Scansiona, si connette, legge e scrive
characteristic & PC Windows con Python + Bleak \\
Peripheral / GATT Server & Espone il servizio GATT PQ-BLE & nRF54L15 DK
con Zephyr \\
Sniffer passivo & Osserva il traffico radio BLE & nRF52840 Dongle con
Wireshark \\
\end{longtable}

\subsection{ML-KEM-768}\label{ml-kem-768}

ML-KEM è lo standard NIST FIPS 203 per key encapsulation post-quantum. È
basato sulla difficoltà del problema Module-LWE. Nel progetto viene
usato il profilo \textbf{ML-KEM-768}, scelto come compromesso tra
sicurezza e overhead.

\begin{longtable}[]{@{}lr@{}}
\toprule\noalign{}
Parametro & Valore \\
\midrule\noalign{}
\endhead
\bottomrule\noalign{}
\endlastfoot
Public key & 1184 byte \\
Secret key & 2400 byte \\
Ciphertext & 1088 byte \\
Shared secret & 32 byte \\
Categoria NIST & 3 \\
\end{longtable}

Un KEM non è uguale a ECDH: il peripheral espone una public key, il
central incapsula un segreto contro quella chiave pubblica e produce un
ciphertext. Chi possiede la secret key può decapsulare il ciphertext e
ottenere la stessa shared secret.

\subsection{SAS Numeric Comparison}\label{sas-numeric-comparison}

Il SAS è un codice breve derivato dalla trascrizione dell'handshake. Nel
progetto è un numero a 6 cifre:

\begin{lstlisting}
transcript = public_key || ciphertext || shared_secret
sas = SHA256(transcript)[0:4] mod 1_000_000
\end{lstlisting}

Se un attaccante MITM sostituisce la public key, central e peripheral
derivano trascrizioni diverse e quindi SAS diversi. Se l'utente
confronta correttamente il codice, la probabilità di falsa accettazione
è circa 1 su 1.000.000.

\subsection{HKDF e AES-256-GCM}\label{hkdf-e-aes-256-gcm}

La shared secret ML-KEM viene trasformata in una chiave di sessione
tramite HKDF-SHA256. I dati applicativi sono poi cifrati con
AES-256-GCM. Nel progetto, l'AAD include:

\begin{lstlisting}
session_id || sender_role || seq_num || msg_type
\end{lstlisting}

Questa scelta lega ogni messaggio alla sessione, alla direzione, al
numero di sequenza e al tipo di messaggio.

\begin{center}\rule{0.5\linewidth}{0.5pt}\end{center}

\section{Modello di minaccia}\label{capitolo-3---modello-di-minaccia}

\subsection{Assunzioni}\label{assunzioni}

Il modello di minaccia considera un attaccante capace di:

\begin{itemize}
\tightlist
\item
  sniffare il traffico BLE;
\item
  registrare pacchetti per un attacco futuro;
\item
  modificare, ritardare o iniettare messaggi GATT;
\item
  tentare un MITM durante l'handshake.
\end{itemize}

Si assume invece che:

\begin{itemize}
\tightlist
\item
  gli endpoint non siano compromessi;
\item
  l'utente confronti correttamente il SAS;
\item
  il generatore casuale del sistema operativo sia sicuro;
\item
  gli attacchi side-channel fisici siano fuori scope.
\end{itemize}

\subsection{Minacce considerate}\label{minacce-considerate}

\begin{longtable}[]{@{}ll@{}}
\toprule\noalign{}
Minaccia & Mitigazione \\
\midrule\noalign{}
\endhead
\bottomrule\noalign{}
\endlastfoot
Store-now-decrypt-later & ML-KEM-768 \\
MITM con sostituzione public key & SAS Numeric Comparison \\
Replay di messaggi cifrati & Sequence number monotono e AAD \\
Tampering del payload & AES-256-GCM authentication tag \\
Cross-session confusion & \passthrough{\lstinline!session\_id!} in
AAD \\
Confusione direzione/messaggio & \passthrough{\lstinline!sender\_role!}
e \passthrough{\lstinline!msg\_type!} in AAD \\
\end{longtable}

\subsection{Minacce fuori scope}\label{minacce-fuori-scope}

\begin{longtable}[]{@{}
  >{\raggedright\arraybackslash}p{(\columnwidth - 2\tabcolsep) * \real{0.5000}}
  >{\raggedright\arraybackslash}p{(\columnwidth - 2\tabcolsep) * \real{0.5000}}@{}}
\toprule\noalign{}
\begin{minipage}[b]{\linewidth}\raggedright
Minaccia
\end{minipage} & \begin{minipage}[b]{\linewidth}\raggedright
Motivazione
\end{minipage} \\
\midrule\noalign{}
\endhead
\bottomrule\noalign{}
\endlastfoot
Jamming BLE & Attacco fisico/radio non mitigabile a livello
applicativo \\
Side-channel su microcontrollore & Richiede contromisure hardware
specifiche \\
Compromissione endpoint & Se l'endpoint è compromesso, la chiave può
essere esfiltrata \\
Errori utente nel SAS & Il protocollo richiede confronto umano
corretto \\
\end{longtable}

\begin{center}\rule{0.5\linewidth}{0.5pt}\end{center}

\section{Protocollo PQ-BLE-HANDSHAKE}\label{capitolo-4---protocollo-pq-ble-handshake}

\subsection{Service GATT}\label{service-gatt}

Il firmware espone un servizio custom:

\begin{lstlisting}
Service UUID: 12345678-1234-1234-1234-123456789abc
\end{lstlisting}

\begin{longtable}[]{@{}
  >{\raggedright\arraybackslash}p{(\columnwidth - 6\tabcolsep) * \real{0.2500}}
  >{\raggedright\arraybackslash}p{(\columnwidth - 6\tabcolsep) * \real{0.2500}}
  >{\raggedright\arraybackslash}p{(\columnwidth - 6\tabcolsep) * \real{0.2500}}
  >{\raggedright\arraybackslash}p{(\columnwidth - 6\tabcolsep) * \real{0.2500}}@{}}
\toprule\noalign{}
\begin{minipage}[b]{\linewidth}\raggedright
Characteristic
\end{minipage} & \begin{minipage}[b]{\linewidth}\raggedright
UUID
\end{minipage} & \begin{minipage}[b]{\linewidth}\raggedright
Proprietà
\end{minipage} & \begin{minipage}[b]{\linewidth}\raggedright
Funzione
\end{minipage} \\
\midrule\noalign{}
\endhead
\bottomrule\noalign{}
\endlastfoot
Public Key & \passthrough{\lstinline!...9abd!} & READ & Espone una public key ML-KEM-768 valida (value handle \texttt{0x0012}) \\
Ciphertext & \passthrough{\lstinline!...9abe!} & WRITE & Riceve il ciphertext ML-KEM dal central (value handle \texttt{0x0014}) \\
Secure Data & \passthrough{\lstinline!...9abf!} & NOTIFY & Invia la notification raw demo (value handle \texttt{0x0016}, CCCD \texttt{0x0017}) \\
Control & \passthrough{\lstinline!...9ac0!} & WRITE & Riceve comandi come \passthrough{\lstinline!START!} (value handle \texttt{0x0019}) \\
\end{longtable}

\subsection{Flusso ideale del
protocollo}\label{flusso-ideale-del-protocollo}

Nel protocollo completo ideale:

\begin{lstlisting}
Peripheral                           Central

keygen() -> pk, sk
espone pk via GATT READ       ->      legge pk

                                  encapsulate(pk) -> ct, ss
riceve ct via GATT WRITE      <-      scrive ct

decapsulate(sk, ct) -> ss

derive_sas(pk, ct, ss)              derive_sas(pk, ct, ss)
confronto umano SAS

HKDF(ss) -> session_key             HKDF(ss) -> session_key

AES-256-GCM applicativo su GATT NOTIFY/WRITE
\end{lstlisting}

\subsection{Stato reale dell'implementazione
hardware}\label{stato-reale-dellimplementazione-hardware}

L'implementazione firmware attuale non esegue ancora il flusso completo
sul DK. Il DK:

\begin{itemize}
\tightlist
\item
  espone il servizio GATT custom e una public key ML-KEM-768 valida da
  1184 byte, generata offline dallo script
  \texttt{scripts/generate\_firmware\_public\_key.py};
\item
  permette al central di verificare il fingerprint SHA-256 della chiave
  letta via BLE, pari a
  \texttt{660da9e73485e45ac35a3a6cd4d65fe39efe161bb01444893733c102566ee912};
\item
  riceve e riassembla il ciphertext da 1088 byte;
\item
  gestisce il comando \passthrough{\lstinline!START!};
\item
  invia una notification raw dimostrativa da 57 byte.
\end{itemize}

La secret key associata alla public key non è stata persistita e non è
presente nella repository. Di conseguenza, la chiave pubblica usata dal
central è autenticamente una chiave ML-KEM, ma il DK non può ancora
eseguire la decapsulazione. Questa scelta evita di presentare come
end-to-end una prova che, allo stato attuale, è deliberatamente limitata
al trasporto.

Il central Python esegue invece:

\begin{itemize}
\tightlist
\item
  ML-KEM encapsulation;
\item
  SAS derivation;
\item
  HKDF;
\item
  AES-GCM secure channel nei test Python.
\end{itemize}

Questo significa che la demo hardware valida il \textbf{trasporto
BLE/GATT reale}, non ancora la crittografia end-to-end completamente
eseguita sul microcontrollore.

\subsection{Frammentazione}\label{frammentazione}

Il materiale ML-KEM è più grande dei payload ATT normalmente
trasportabili in un singolo pacchetto. Il progetto usa un header
applicativo da 4 byte:

\begin{lstlisting}
idx (1 byte) || total (1 byte) || payload_len (2 byte)
\end{lstlisting}

Con MTU 247, validato nella demo Windows, il payload utile per frammento
ciphertext è:

\begin{lstlisting}
247 - 4 = 243 byte
\end{lstlisting}

Quindi il ciphertext ML-KEM da 1088 byte viene scritto in:

\begin{lstlisting}
ceil(1088 / 243) = 5 frammenti
\end{lstlisting}

La public key da 1184 byte viene invece letta tramite long read ATT, con
read blob successivi.

\begin{center}\rule{0.5\linewidth}{0.5pt}\end{center}

\section{Implementazione}\label{capitolo-5---implementazione}

\subsection{Stack tecnologico}\label{stack-tecnologico}

\begin{longtable}[]{@{}
  >{\raggedright\arraybackslash}p{(\columnwidth - 4\tabcolsep) * \real{0.3333}}
  >{\raggedright\arraybackslash}p{(\columnwidth - 4\tabcolsep) * \real{0.3333}}
  >{\raggedright\arraybackslash}p{(\columnwidth - 4\tabcolsep) * \real{0.3333}}@{}}
\toprule\noalign{}
\begin{minipage}[b]{\linewidth}\raggedright
Componente
\end{minipage} & \begin{minipage}[b]{\linewidth}\raggedright
Tecnologia
\end{minipage} & \begin{minipage}[b]{\linewidth}\raggedright
Ruolo
\end{minipage} \\
\midrule\noalign{}
\endhead
\bottomrule\noalign{}
\endlastfoot
Python & 3.13 su Windows, 3.x su Linux & Implementazione protocollo e
test \\
liboqs-python / liboqs & ML-KEM-768 & Primitive post-quantum \\
cryptography & HKDF, AES-GCM & Canale cifrato applicativo \\
bleak & BLE central & Scan, connect, read/write/notify \\
Zephyr / nRF Connect SDK & Firmware nRF54L15 DK & Peripheral GATT \\
Wireshark + nRF Sniffer & Cattura passiva & Validazione
osservazionale \\
\end{longtable}

\subsection{Struttura del
repository}\label{struttura-del-repository}

\begin{lstlisting}
pq-ble-handshake/
|--- src/
|   |--- common/                  # ML-KEM, SAS, sessioni, frammentazione
|   `--- central/                 # Bleak central e demo hardware
|--- firmware/
|   `--- nrf54l15_pq_gatt_skeleton/
|       |--- src/main.c
|       `--- src/demo_public_key.h
|--- scripts/
|   `--- generate_firmware_public_key.py
|--- tests/                       # 101 test attivi
|--- benchmarks/
|   |--- run_all.ps1
|   `--- results/                 # output TXT e JSON misurati
|--- docs/
|   |--- images/                  # cinque screenshot Wireshark
|   |--- captures/                # catture PCAPNG
|   `--- hardware-validation-log.txt
|--- report/                      # sorgente LaTeX e PDF
`--- data/keys/                   # store sessioni locale, non versionato
\end{lstlisting}

La cartella \passthrough{\lstinline!experimental/peripheral/!} non è
parte della demo hardware reale. La demo reale usa il firmware Zephyr su
nRF54L15 DK.

\subsection{Moduli Python
principali}\label{moduli-python-principali}

\subsubsection{\texorpdfstring{\texttt{common/ml\_kem.py}}{common/ml\_kem.py}}\label{commonml_kem.py}

Wrapper per ML-KEM-768. Le funzioni principali sono:

\begin{lstlisting}[language=Python]
generate_keypair() -> (public_key, secret_key)
encapsulate(public_key) -> (ciphertext, shared_secret)
decapsulate(secret_key, ciphertext) -> shared_secret
\end{lstlisting}

Durante la validazione su Windows è stata aggiunta la conversione
esplicita \passthrough{\lstinline!bytes(...)!} sugli input e sugli
output, perché Bleak può restituire dati GATT come
\passthrough{\lstinline!bytearray!}, mentre
\passthrough{\lstinline!liboqs-python!} si aspetta buffer compatibili
con \passthrough{\lstinline!bytes!}.

\begin{lstlisting}[style=pythonstyle,caption={Normalizzazione dei buffer per liboqs.}]
public_key = bytes(public_key)
ciphertext = bytes(ciphertext)
\end{lstlisting}

\subsubsection{\texorpdfstring{\texttt{common/fragmentation.py}}{common/fragmentation.py}}\label{commonfragmentation.py}

Implementa:

\begin{itemize}
\tightlist
\item
  \passthrough{\lstinline!Fragment.encode()!};
\item
  \passthrough{\lstinline!Fragment.decode()!};
\item
  \passthrough{\lstinline!fragment\_data()!};
\item
  \passthrough{\lstinline!reassemble\_data()!}.
\end{itemize}

Gestisce casi di errore come frammenti mancanti, duplicati o header non
validi. Il formato applicativo è esplicito e indipendente dal mapping
interno scelto dal backend ATT:

\begin{lstlisting}[style=pythonstyle,caption={Header applicativo dei frammenti PQ-BLE.}]
header = bytes([index, total]) + payload_len.to_bytes(2, "big")
wire_fragment = header + payload
\end{lstlisting}

\subsubsection{\texorpdfstring{\texttt{common/sas.py}}{common/sas.py}}\label{commonsas.py}

Deriva e formatta il SAS a 6 cifre. I test verificano determinismo,
sensibilità ai cambiamenti e formato.

\begin{lstlisting}[style=pythonstyle,caption={Derivazione del SAS dalla trascrizione.}]
transcript = public_key + ciphertext + shared_secret
digest = hashlib.sha256(transcript).digest()
sas = int.from_bytes(digest[:4], "big") % 1_000_000
\end{lstlisting}

\subsubsection{\texorpdfstring{\texttt{common/session.py}}{common/session.py}}\label{commonsession.py}

Implementa:

\begin{itemize}
\tightlist
\item
  derivazione della session key con HKDF-SHA256;
\item
  canale AES-256-GCM;
\item
  AAD con \passthrough{\lstinline!session\_id!}, ruolo, sequence number
  e tipo messaggio;
\item
  replay protection e rifiuto di messaggi fuori ordine.
\end{itemize}

\begin{lstlisting}[style=pythonstyle,caption={Costruzione dell'Additional Authenticated Data.}]
def _build_aad(self, seq_num, sender_role, msg_type):
    return (
        self.session_id
        + sender_role
        + seq_num.to_bytes(8, "big")
        + msg_type
    )
\end{lstlisting}

Il wire format risultante aggiunge 37 byte al plaintext: 8 byte di
sequence number, 1 byte di tipo, 12 byte di IV e 16 byte di tag GCM.

\subsubsection{\texorpdfstring{\texttt{common/session\_store.py}}{common/session\_store.py}}\label{commonsession_store.py}

Implementa uno store JSON per session resumption. Ogni sessione salva:

\begin{itemize}
\tightlist
\item
  \passthrough{\lstinline!session\_id!};
\item
  \passthrough{\lstinline!session\_key!};
\item
  timestamp di creazione;
\item
  contatore di resume riusciti.
\end{itemize}

Una sessione è invalidata al raggiungimento del primo dei due limiti:
24 ore oppure 100 resume riusciti. La lettura preliminare non incrementa
il contatore; l'incremento avviene soltanto dopo una risposta
\texttt{RESUME\_OK}. Timeout e risposte negative non consumano un uso.

\begin{lstlisting}[style=pythonstyle,caption={Il contatore viene consumato solo dopo un resume accettato.}]
session_key = store.load(session_id, increment_usage=False)
response = await wait_resume_response()
if response == RESUME_OK_NOTIFY:
    store.mark_used(session_id)
\end{lstlisting}

\subsubsection{\texorpdfstring{\texttt{central/ble\_client.py}}{central/ble\_client.py}}\label{centralble_client.py}

Gestisce il central BLE:

\begin{itemize}
\tightlist
\item
  scan del dispositivo;
\item
  connessione;
\item
  lettura della public key;
\item
  scrittura frammentata del ciphertext;
\item
  invio comandi di controllo;
\item
  subscription alle notification.
\end{itemize}

\subsubsection{\texorpdfstring{\texttt{central/handshake.py}}{central/handshake.py}}\label{centralhandshake.py}

Orchestra il protocollo lato central:

\begin{enumerate}
\def\labelenumi{\arabic{enumi}.}
\tightlist
\item
  tentativo di session resume;
\item
  fallback a full handshake;
\item
  lettura public key;
\item
  ML-KEM encapsulation;
\item
  scrittura ciphertext;
\item
  SAS;
\item
  derivazione session key;
\item
  salvataggio sessione.
\end{enumerate}

\subsubsection{\texorpdfstring{\texttt{central/secure\_channel.py}}{central/secure\_channel.py}}\label{centralsecure_channel.py}

Gestisce la ricezione delle notification. È stata aggiunta una modalità:

\begin{lstlisting}[language=Python]
start_receiving(decrypt_notifications=False)
\end{lstlisting}

Questa modalità è usata nella demo hardware, perché il firmware attuale
invia una notification raw e non un payload AES-GCM coerente con la
session key del central.

\subsection{Firmware nRF54L15 DK}\label{firmware-nrf54l15-dk}

Il firmware Zephyr in
\passthrough{\lstinline!firmware/nrf54l15\_pq\_gatt\_skeleton/!}
implementa un peripheral BLE reale.

Funzionalità validate:

\begin{itemize}
\tightlist
\item
  advertising come \passthrough{\lstinline!PQ-BLE-Device!};
\item
  service custom PQ-BLE;
\item
  Public Key READ;
\item
  Ciphertext WRITE;
\item
  Control WRITE;
\item
  Secure Data NOTIFY;
\item
  CCCD;
\item
  gestione MTU;
\item
  accumulo frammenti ciphertext;
\item
  invio reale con \passthrough{\lstinline!bt\_gatt\_notify()!}.
\end{itemize}

Configurazione rilevante:

\begin{lstlisting}
CONFIG_BT=y
CONFIG_BT_PERIPHERAL=y
CONFIG_BT_DEVICE_NAME="PQ-BLE-Device"
CONFIG_BT_SMP=n
\end{lstlisting}

\passthrough{\lstinline!CONFIG\_BT\_SMP=n!} è intenzionale: il progetto
dimostra una protezione applicativa sopra GATT e non usa il pairing
BLE/SMP.

\subsection{Stato del firmware}\label{stato-del-firmware}

Il firmware è stato compilato e flashato con successo su nRF54L15 DK.
La versione corrente espone una public key ML-KEM-768 valida generata
offline e ne rende verificabile il fingerprint attraverso il log del
central. Il materiale segreto non viene salvato. Il firmware resta quindi
una \textbf{BLE/GATT transport validation firmware}: riceve e riassembla
il ciphertext, ma non esegue ancora:

\begin{itemize}
\tightlist
\item
  key generation o decapsulation ML-KEM on-chip;
\item
  HKDF e AES-GCM on-chip;
\item
  session store persistente su flash;
\item
  generazione di una notification cifrata con la stessa session key del
  central.
\end{itemize}

La notification da 57 byte è esplicitamente un payload raw dimostrativo.
Questa distinzione impedisce di sovrastimare il risultato sperimentale.

\begin{center}\rule{0.5\linewidth}{0.5pt}\end{center}

\section{Validazione sperimentale e benchmark}\label{capitolo-6---validazione-sperimentale-e-benchmark}

\subsection{Test automatizzati}\label{test-automatizzati}

La suite attiva contiene \textbf{101 test passanti}.

\begin{lstlisting}[language=bash]
python -m pytest tests/ -v
\end{lstlisting}

\begin{longtable}[]{@{}
  >{\raggedright\arraybackslash}p{(\columnwidth - 4\tabcolsep) * \real{0.3000}}
  >{\raggedleft\arraybackslash}p{(\columnwidth - 4\tabcolsep) * \real{0.4000}}
  >{\raggedright\arraybackslash}p{(\columnwidth - 4\tabcolsep) * \real{0.3000}}@{}}
\toprule\noalign{}
\begin{minipage}[b]{\linewidth}\raggedright
File
\end{minipage} & \begin{minipage}[b]{\linewidth}\raggedleft
Test attivi
\end{minipage} & \begin{minipage}[b]{\linewidth}\raggedright
Cosa verifica
\end{minipage} \\
\midrule\noalign{}
\endhead
\bottomrule\noalign{}
\endlastfoot
\passthrough{\lstinline!test\_ml\_kem.py!} & 6 & ML-KEM keygen,
encapsulation, decapsulation, roundtrip \\
\passthrough{\lstinline!test\_fragmentation.py!} & 14 & Frammentazione,
riassemblaggio, missing, duplicati, MTU variabili \\
\passthrough{\lstinline!test\_sas.py!} & 12 & SAS a 6 cifre,
determinismo, sensibilità, distribuzione \\
\passthrough{\lstinline!test\_session.py!} & 22 & HKDF, AES-GCM, AAD,
replay, out-of-order, tampering \\
\passthrough{\lstinline!test\_session\_store.py!} & 23 & Session ID,
resume, store, expiry, usage counter \\
\passthrough{\lstinline!test\_handshake\_mock.py!} & 2 & Pipeline
completa senza BLE reale \\
\passthrough{\lstinline!test\_mitm\_simulation.py!} & 2 & MITM detection
tramite SAS mismatch \\
\passthrough{\lstinline!test\_firmware\_uuids.py!} & 3 & Device name,
SMP disabled, \passthrough{\lstinline!bt\_gatt\_notify()!} \\
\passthrough{\lstinline!test\_central\_transport\_mock.py!} & 17 &
Read/write frammentati con mock GATT \\
\textbf{Totale} & \textbf{101} & Validazione proof-of-concept \\
\end{longtable}

Alcuni test legacy di parsing automatico degli UUID firmware sono stati
commentati perché dipendevano da un formato precedente delle
dichiarazioni C. La coerenza degli UUID è verificata tramite sorgente
firmware, nRF Connect Mobile e cattura reale.

\subsection{Test MITM}\label{test-mitm}

La simulazione MITM genera una keypair legittima e una keypair
dell'attaccante. Il central incapsula contro la public key sostituita
dal MITM, mentre il peripheral legittimo usa una trascrizione diversa.
Il risultato atteso è un SAS diverso.

\begin{lstlisting}
pk_A, sk_A = keypair legittima
pk_M, sk_M = keypair MITM

Central usa pk_M
Peripheral legittimo usa pk_A

sas_central != sas_peripheral
\end{lstlisting}

Questo dimostra che un MITM che sostituisce la public key viene rilevato
se il SAS viene confrontato correttamente.

\subsection{Validazione firmware
build/flash}\label{validazione-firmware-buildflash}

Il firmware è stato compilato su Windows usando nRF Connect SDK e
target:

\begin{lstlisting}
nrf54l15dk/nrf54l15/cpuapp
\end{lstlisting}

Per evitare problemi di path length su Windows, la build è stata
eseguita anche da:

\begin{lstlisting}
C:/myfw/pq
\end{lstlisting}

Risultato della validazione:

\begin{lstlisting}
Firmware build: OK
Firmware flash: OK
Board flashed successfully
\end{lstlisting}

\subsection{Validazione manuale con nRF Connect
Mobile}\label{validazione-manuale-con-nrf-connect-mobile}

Da telefono è stato verificato che:

\begin{itemize}
\tightlist
\item
  il device appare come \passthrough{\lstinline!PQ-BLE-Device!};
\item
  la connessione BLE avviene correttamente;
\item
  la public key è leggibile;
\item
  le notification sono abilitabili;
\item
  il comando \passthrough{\lstinline!START!} viene ricevuto.
\end{itemize}

Log rappresentativo lato DK:

\begin{lstlisting}
PK read: offset=0, len=246
PK read: offset=246, len=246
PK read: offset=492, len=246
Notifications ENABLED
Control write: len=5
Control data: START
Received START command
START received but ciphertext not yet written
\end{lstlisting}

L'ultimo messaggio è corretto se \passthrough{\lstinline!START!} viene
inviato prima del ciphertext.

\subsection{Demo hardware PC central \textless-\textgreater{}
nRF54L15 DK}\label{demo-hardware-pc-central---nrf54l15-dk}

La demo reale è stata eseguita con:

\begin{lstlisting}[language=bash]
python -m src.central.main --device PQ-BLE-Device --demo --no-sas-confirm --log-level DEBUG
\end{lstlisting}

Risultato ottenuto:

\begin{lstlisting}
Found PQ-BLE-Device
Connected
Read public key: 1184 bytes
Public key SHA-256: 660da9e73485e45ac35a3a6cd4d65fe39efe161bb01444893733c102566ee912
Encapsulate: ct=1088 bytes, ss=32 bytes
Writing ciphertext: 1088 bytes in 5 fragments
Ciphertext written
SAS derived
Session key: 32 bytes
START sent
Raw demo notification received: 57 bytes
BLE/GATT transport validation completed.
\end{lstlisting}

Questa prova dimostra che:

\begin{enumerate}
\def\labelenumi{\arabic{enumi}.}
\tightlist
\item
  il PC scopre il DK;
\item
  la connessione BLE avviene;
\item
  la public key ML-KEM-768 valida da 1184 byte viene letta e il suo fingerprint SHA-256 viene verificato;
\item
  il central genera un ciphertext ML-KEM da 1088 byte;
\item
  il ciphertext viene scritto al DK in 5 frammenti;
\item
  il comando \passthrough{\lstinline!START!} viene inviato;
\item
  il DK risponde con una notification raw da 57 byte.
\end{enumerate}

\subsection{Cattura Wireshark con nRF52840
Dongle}\label{cattura-wireshark-con-nrf52840-dongle}

È stata acquisita una cattura passiva con nRF52840 Dongle e Wireshark,
filtrando i pacchetti ATT/GATT con:

\begin{lstlisting}
btatt
\end{lstlisting}

Nella cattura si osservano:

\begin{itemize}
\tightlist
\item
  MTU exchange;
\item
  lettura della public key tramite ATT Read / Read Blob;
\item
  trasferimento del ciphertext sulla characteristic dedicata;
\item
  write del comando \passthrough{\lstinline!START!};
\item
  Handle Value Notification finale.
\end{itemize}

\begin{figure}[H]
  \centering
  \safeincludegraphics[width=0.96\textwidth]{wireshark_01_mtu_exchange.png}
  \caption{Negoziazione dell'ATT MTU tra il central Windows e il peripheral nRF54L15 DK. Il valore negoziato nella demo è 247 byte.}
  \label{fig:wireshark-mtu}
\end{figure}

Handle osservati nella cattura:

\begin{longtable}[]{@{}ll@{}}
\toprule\noalign{}
Funzione & Handle \\
\midrule\noalign{}
\endhead
\bottomrule\noalign{}
\endlastfoot
Public Key characteristic value & \passthrough{\lstinline!0x0012!} \\
Ciphertext characteristic value & \passthrough{\lstinline!0x0014!} \\
Secure Data characteristic value / notification & \passthrough{\lstinline!0x0016!} \\
CCCD notification & \passthrough{\lstinline!0x0017!} \\
Control / START & \passthrough{\lstinline!0x0019!} \\
\end{longtable}

La public key da 1184 byte viene letta con long read ATT. Con MTU 247,
ogni read blob trasporta circa 246 byte utili, per cui si osservano
offset successivi come:

\begin{lstlisting}
0
246
492
738
984
\end{lstlisting}

\begin{figure}[H]
  \centering
  \safeincludegraphics[width=0.96\textwidth]{wireshark_02_public_key_read.png}
  \caption{Lettura lunga del materiale pubblico ML-KEM sulla characteristic con handle \texttt{0x0012}. Dopo la Read Request iniziale, le Read Blob Request usano offset crescenti.}
  \label{fig:wireshark-public-key}
\end{figure}

Il ciphertext da 1088 byte viene scritto sulla characteristic
\passthrough{\lstinline!0x0014!}. Su Windows/Bleak la scrittura lunga
può apparire in Wireshark come sequenza di
\passthrough{\lstinline!Prepare Write Request!},
\passthrough{\lstinline!Prepare Write Response!},
\passthrough{\lstinline!Execute Write Request!} ed
\passthrough{\lstinline!Execute Write Response!}, oltre alla
frammentazione L2CAP/ATT interna. Questo non contraddice la
frammentazione applicativa: il codice PQ-BLE divide il ciphertext in 5
frammenti logici, mentre lo stack BLE può ulteriormente mappare queste
scritture in procedure ATT.

\begin{figure}[H]
  \centering
  \safeincludegraphics[width=0.96\textwidth]{wireshark_03_ciphertext_write.png}
  \caption{Trasferimento del ciphertext ML-KEM sulla characteristic con handle \texttt{0x0014}. Lo stack Windows/Bleak rappresenta le scritture lunghe mediante Prepare Write ed Execute Write.}
  \label{fig:wireshark-ciphertext}
\end{figure}

Il comando \passthrough{\lstinline!START!} è visibile come payload
ASCII:

\begin{lstlisting}
53 54 41 52 54
\end{lstlisting}

\begin{figure}[H]
  \centering
  \safeincludegraphics[width=0.96\textwidth]{wireshark_04_start_write.png}
  \caption{Write Request del comando \texttt{START} sull'handle \texttt{0x0019}. Il valore esadecimale \texttt{5354415254} corrisponde alla stringa ASCII.}
  \label{fig:wireshark-start}
\end{figure}

\begin{figure}[H]
  \centering
  \safeincludegraphics[width=0.96\textwidth]{wireshark_05_notification_57_bytes.png}
  \caption{Handle Value Notification da 57 byte inviata dal DK sull'handle valore \texttt{0x0016}.}
  \label{fig:wireshark-notification}
\end{figure}

La notification finale è visibile come ATT Handle Value Notification. Il value handle è \texttt{0x0016}, mentre \texttt{0x0017} è il CCCD usato per abilitare le notifiche.

\subsection{Interpretazione della
cattura}\label{interpretazione-della-cattura}

La cattura non dimostra che il DK esegua ML-KEM o AES-GCM on-chip.
Dimostra invece che il protocollo viene effettivamente trasportato su
BLE ATT/GATT standard:

\begin{lstlisting}
Public key read        -> ATT Read / Read Blob
Ciphertext transfer    -> ATT Write / Prepare Write / Execute Write
Control START          -> ATT Write
Firmware response      -> ATT Handle Value Notification
\end{lstlisting}

Questo è esattamente l'obiettivo della validazione hardware attuale.

\subsection{Benchmark riproducibili}\label{benchmark}

La suite in \texttt{benchmarks/} distingue deliberatamente il costo CPU
delle primitive dal tempo end-to-end BLE. I risultati vengono prodotti in
formato TXT e JSON insieme alla descrizione dell'ambiente di esecuzione
(CPU, sistema operativo, versione Python, dipendenze e commit Git). Il
comando Windows è:

\begin{lstlisting}[style=bashstyle]
.\benchmarks\run_all.ps1
\end{lstlisting}

\ifbenchmarkdata
I valori seguenti provengono dagli artefatti JSON reali generati il
\BenchGeneratedAt. Il benchmark crittografico ha utilizzato
\BenchIterations{} iterazioni misurate, escludendo warm-up, scan,
connessione e trasferimento GATT.
\else
\begin{importantbox}
I file JSON dei benchmark non erano disponibili durante la generazione di
questa copia. La struttura è pronta, ma i valori CPU restano marcati come
\texttt{N/D}. Per popolarli senza modifiche manuali eseguire
\texttt{python scripts/generate\_benchmark\_latex.py} dalla root della
repository e ricompilare il documento.
\end{importantbox}
\fi

\subsubsection{Latenza crittografica del full handshake}

\begin{table}[H]
\centering
\caption{Latenza CPU delle primitive del full handshake. I tempi sono in microsecondi, salvo il totale espresso in millisecondi.}
\label{tab:benchmark-handshake}
\begin{tabular}{lrrr}
\toprule
Operazione & Media & Mediana & P99 \\
\midrule
ML-KEM key generation & \BenchKeygenMean & \BenchKeygenMedian & \BenchKeygenPnn \\
ML-KEM encapsulation & \BenchEncapsMean & \BenchEncapsMedian & \BenchEncapsPnn \\
ML-KEM decapsulation & \BenchDecapsMean & \BenchDecapsMedian & \BenchDecapsPnn \\
SAS derivation & \BenchSasMean & -- & -- \\
HKDF-SHA256 & \BenchHkdfMean & -- & -- \\
\midrule
Totale crittografico [ms] & \BenchTotalMeanMs & \BenchTotalMedianMs & \BenchTotalPnnMs \\
\bottomrule
\end{tabular}
\end{table}

Il throughput derivato dalla media è pari a
\BenchHandshakeRate{} full handshake crittografici al secondo sulla
macchina di prova. Questo valore non è una misura del numero di handshake
BLE realizzabili al secondo: nel percorso reale dominano anche discovery,
connessione, procedure ATT e scheduling del sistema operativo.

\subsubsection{Throughput CPU del canale AES-256-GCM}

\begin{table}[H]
\centering
\caption{Throughput CPU del SecureChannel. I valori sono in KiB/s e non rappresentano il throughput radio BLE.}
\label{tab:benchmark-throughput}
\begin{tabular}{rrrr}
\toprule
Payload [B] & Cifratura & Decifratura & Overhead wire \\
\midrule
64 & \BenchEncSixtyFour & \BenchDecSixtyFour & 57,81\% \\
256 & \BenchEncTwoFiftySix & \BenchDecTwoFiftySix & 14,45\% \\
512 & \BenchEncFiveTwelve & \BenchDecFiveTwelve & 7,23\% \\
1024 & \BenchEncOneK & \BenchDecOneK & 3,61\% \\
4096 & \BenchEncFourK & \BenchDecFourK & 0,90\% \\
16384 & \BenchEncSixteenK & \BenchDecSixteenK & 0,23\% \\
\bottomrule
\end{tabular}
\end{table}

L'overhead fisso è di 37 byte per messaggio:

\begin{equation}
8_{\mathrm{seq}} + 1_{\mathrm{type}} + 12_{\mathrm{IV}} + 16_{\mathrm{tag}} = 37\ \mathrm{byte}.
\end{equation}

Per payload piccoli l'incidenza percentuale è elevata; cresce invece
l'efficienza all'aumentare della dimensione del plaintext.

\subsubsection{Frammentazione applicativa}

\begin{table}[H]
\centering
\caption{Overhead del ciphertext ML-KEM-768 nella frammentazione applicativa.}
\label{tab:benchmark-fragmentation}
\begin{tabular}{rrrrrr}
\toprule
MTU & Payload/frag. & Frammenti & Wire [B] & Overhead & Frag./reasm. [\textmu s] \\
\midrule
247 & 243 & 5 & 1108 & 20 B (1,84\%) & \BenchFragMtuTwoFourSevenUs{} / \BenchReasmMtuTwoFourSevenUs \\
512 & 508 & 3 & 1100 & 12 B (1,10\%) & \BenchFragMtuFiveTwelveUs{} / \BenchReasmMtuFiveTwelveUs \\
\bottomrule
\end{tabular}
\end{table}

Il caso MTU 247 corrisponde alla demo reale. La public key non percorre
questa funzione di frammentazione: viene letta mediante ATT Long Read e
Read Blob, come mostrato dalla cattura Wireshark. I tempi di
frammentazione e riassemblaggio misurano solamente il codice Python; non
includono la latenza del link radio.
\subsubsection{Distribuzione dell'overhead tra handshake e comunicazione}

È importante distinguere il costo del \textit{key establishment}
post-quantum dal costo della successiva comunicazione applicativa.
ML-KEM-768 non viene utilizzato per cifrare ogni messaggio BLE, ma
soltanto per permettere ai due endpoint di ottenere una stessa
\textit{shared secret} senza trasmetterla direttamente sul canale.

Durante un full handshake vengono trasferiti i due principali oggetti
ML-KEM-768:

\[
B_{\mathrm{KEM}}
=
B_{\mathrm{pk}} + B_{\mathrm{ct}}
=
1184 + 1088
=
2272 \text{ byte}.
\]

Questo valore rappresenta il materiale crittografico principale e non
include gli overhead dei livelli ATT, L2CAP e Link Layer.

Nel caso realmente validato con MTU 247, il ciphertext viene suddiviso
in cinque frammenti applicativi. Poiché ogni frammento contiene un header
da 4 byte, il costo applicativo del trasferimento diventa:

\[
B_{\mathrm{full}}
=
1184 + 1088 + (5 \cdot 4)
=
2292 \text{ byte}.
\]

La public key non utilizza la frammentazione applicativa implementata dal
progetto, ma viene trasferita mediante ATT Long Read e successive
Read Blob Request. Il valore di 2292 byte non rappresenta quindi il
numero complessivo di byte trasmessi via radio, ma il materiale
applicativo trasportato da PQ-BLE-HANDSHAKE.

Una volta ottenuta la shared secret, entrambi gli endpoint derivano una
session key da 32 byte tramite HKDF-SHA256:

\[
K_{\mathrm{session}}
=
\operatorname{HKDF\text{-}SHA256}
\left(
SS_{\mathrm{ML\text{-}KEM}}
\right).
\]

ML-KEM non viene più utilizzato durante il trasferimento ordinario dei
dati. Ogni messaggio applicativo viene invece protetto mediante
AES-256-GCM. Nel wire format implementato dal proof-of-concept vengono
aggiunti:

\[
8_{\mathrm{seq}}
+
1_{\mathrm{type}}
+
12_{\mathrm{IV}}
+
16_{\mathrm{tag}}
=
37 \text{ byte}.
\]

Per un plaintext di lunghezza $L$, la dimensione del messaggio protetto è
quindi:

\[
B_{\mathrm{secure}}(L) = L + 37.
\]

\begin{table}[H]
    \centering
    \small
    \caption{Distribuzione dell'overhead nelle diverse fasi del protocollo.}
    \label{tab:protocol-overhead-distribution}
    \begin{tabular}{p{3.5cm}p{4.1cm}p{3.2cm}p{3.3cm}}
        \hline
        \textbf{Fase} &
        \textbf{Materiale principale} &
        \textbf{Costo applicativo} &
        \textbf{Primitive utilizzate} \\
        \hline

        Full handshake &
        Public key, ciphertext e header di frammentazione &
        2292 byte &
        ML-KEM-768, SAS, HKDF \\

        Comunicazione nella stessa sessione &
        Plaintext cifrato e autenticato &
        $L + 37$ byte per messaggio &
        AES-256-GCM \\

        Session resumption &
        Resume request e risposta positiva &
        30 byte di controllo, prima degli overhead GATT &
        Chiave cached e AES-256-GCM \\

        Re-handshake dopo scadenza &
        Nuova public key e nuovo ciphertext &
        2292 byte &
        Nuovo ML-KEM-768 \\
        \hline
    \end{tabular}
\end{table}

Nel protocollo Python attuale, una richiesta di resume occupa 21 byte:

\[
4_{\mathrm{magic}}
+
1_{\mathrm{type}}
+
16_{\mathrm{session\_id}}
=
21 \text{ byte},
\]

mentre la risposta \texttt{RESUME\_OK} occupa 9 byte. Il controllo di
resume richiede quindi 30 byte applicativi, escludendo gli overhead
ATT/GATT. Questo valore è molto inferiore ai 2292 byte del full
handshake.

La session resumption consente quindi di concentrare il maggiore costo
post-quantum nel primo handshake e nei successivi re-handshake
periodici. Una connessione ripresa utilizza la chiave già memorizzata e
prosegue con AES-256-GCM, senza trasferire nuovamente la public key e il
ciphertext ML-KEM.

La comunicazione rimane resistente allo scenario
\textit{store-now-decrypt-later}, perché la session key originaria è
stata stabilita mediante ML-KEM-768 e i dati sono protetti con una
chiave AES da 256 bit. La resumption presenta tuttavia un diverso
trade-off di sicurezza: il riutilizzo della chiave memorizzata non offre
forward secrecy completa. Se la session key cached viene sottratta da un
endpoint, può essere compromesso tutto il traffico protetto con quella
chiave fino al successivo full handshake.

Nel proof-of-concept corrente la session resumption completa è validata
soltanto nella componente Python. Il firmware nRF54L15 DK non dispone
ancora di uno store persistente e non esegue la derivazione della
session key o AES-256-GCM on-chip. I 30 byte rappresentano quindi il
costo previsto dal protocollo implementato e testato lato software, non
una misura ottenuta dalla demo hardware.

\begin{table}[H]
    \centering
    \small
    \caption{Incidenza dell'overhead fisso del SecureChannel.}
    \label{tab:secure-channel-overhead}
    \begin{tabular}{r r r}
        \hline
        \textbf{Plaintext} &
        \textbf{Messaggio protetto} &
        \textbf{Overhead} \\
        \hline
        20 B   & 57 B   & 185,00\% \\
        64 B   & 101 B  & 57,81\% \\
        256 B  & 293 B  & 14,45\% \\
        512 B  & 549 B  & 7,23\% \\
        1024 B & 1061 B & 3,61\% \\
        \hline
    \end{tabular}
\end{table}

L'overhead fisso risulta particolarmente significativo per i piccoli
valori tipici dei sensori BLE, mentre diventa progressivamente meno
rilevante per payload di alcune centinaia di byte. Questo costo non
dipende direttamente da ML-KEM, ma dal wire format applicativo scelto
per AES-256-GCM, replay protection e separazione dei messaggi.

\begin{center}\rule{0.5\linewidth}{0.5pt}\end{center}

\section{Discussione e limiti}\label{capitolo-7---discussione-e-limiti}

\subsection{Cosa è stato
dimostrato}\label{cosa-uxe8-stato-dimostrato}

Il progetto dimostra:

\begin{itemize}
\tightlist
\item
  correttezza della logica ML-KEM/SAS/HKDF/AES-GCM in Python;
\item
  replay protection e AAD binding;
\item
  session resumption lato Python con limiti di 24 ore e 100 resume riusciti;
\item
  public key ML-KEM-768 valida nel firmware e verifica del fingerprint via BLE;
\item
  benchmark riproducibili con output TXT/JSON;
\item
  trasferimento reale via BLE/GATT tra PC e nRF54L15 DK;
\item
  catturabilità e osservabilità del traffico con nRF52840/Wireshark.
\end{itemize}

\subsection{Cosa non è stato ancora
dimostrato}\label{cosa-non-uxe8-stato-ancora-dimostrato}

Il progetto non dimostra ancora:

\begin{itemize}
\tightlist
\item
  ML-KEM keygen on-chip sul DK;
\item
  ML-KEM decapsulation on-chip;
\item
  HKDF e AES-GCM on-chip;
\item
  cifratura end-to-end di una notification generata dal DK con la stessa
  session key del central;
\item
  persistenza della session key su flash del DK;
\item
  consumo energetico reale del protocollo embedded completo.
\end{itemize}

\subsection{Session resumption e forward secrecy}\label{session-resumption}

La session resumption riduce il costo delle riconnessioni memorizzando:

\begin{lstlisting}
session_id -> session_key
\end{lstlisting}

Il central presenta il \texttt{session\_id}; solo una risposta
\texttt{RESUME\_OK} incrementa il contatore. Timeout e rifiuti non
consumano un uso. Una sessione viene invalidata dopo 24 ore oppure al
raggiungimento di 100 resume riusciti. Il centesimo resume è accettato,
poi l'entry viene eliminata e la connessione successiva forza un nuovo
full handshake ML-KEM.

Nel firmware corrente il resume non è ancora implementato, perciò i
tentativi hardware terminano per timeout e il central esegue il fallback.
In produzione converrebbe inoltre provare prima la sessione più recente e
usare un timeout breve per evitare ritardi quando sono presenti più entry
cached.

\paragraph{Trade-off di forward secrecy}
Il resume riutilizza una session key persistita e non offre quindi forward
secrecy piena alle sessioni riassunte. Se la chiave cached viene
compromessa, può essere esposto il traffico protetto con quella chiave
fino al successivo full handshake. I limiti temporale e di utilizzo
riducono la finestra di esposizione, ma non equivalgono alla generazione
di materiale effimero nuovo per ogni connessione. Una versione più forte
potrebbe derivare sottochiavi per connessione e cancellare in modo sicuro
lo stato precedente.

È importante distinguere la resistenza quantistica dalla forward secrecy.
Una sessione ripresa rimane protetta contro un attaccante passivo che
abbia registrato il traffico, poiché la chiave cached è stata inizialmente
derivata da ML-KEM-768 e non da ECDH P-256. Un futuro quantum computer
non può quindi applicare l'algoritmo di Shor al traffico dell'handshake
per ricostruire la session key.

La resumption non genera però nuovo materiale post-quantum e riutilizza
la stessa session key. La compromissione dello stato persistente di uno
degli endpoint espone pertanto tutte le connessioni che usano quella
chiave. I limiti di 24 ore e 100 resume riducono questa finestra, ma non
forniscono forward secrecy piena.

\subsection{Confronto con ECDH BLE
standard}\label{confronto-con-ecdh-ble-standard}

\begin{longtable}[]{@{}
  >{\raggedright\arraybackslash}p{(\columnwidth - 4\tabcolsep) * \real{0.3333}}
  >{\raggedright\arraybackslash}p{(\columnwidth - 4\tabcolsep) * \real{0.3333}}
  >{\raggedright\arraybackslash}p{(\columnwidth - 4\tabcolsep) * \real{0.3333}}@{}}
\toprule\noalign{}
\begin{minipage}[b]{\linewidth}\raggedright
Proprietà
\end{minipage} & \begin{minipage}[b]{\linewidth}\raggedright
ECDH-P256 BLE
\end{minipage} & \begin{minipage}[b]{\linewidth}\raggedright
PQ-BLE-HANDSHAKE
\end{minipage} \\
\midrule\noalign{}
\endhead
\bottomrule\noalign{}
\endlastfoot
Resistenza quantistica & No & Sì, lato protocollo ML-KEM \\
Standard Bluetooth & Sì & No, proof-of-concept applicativo \\
Materiale principale di key establishment &
Due public key ECDH P-256, 128 B complessivi &
Public key ML-KEM e ciphertext, 2272 B \\

Cifratura dei dati successivi &
AES-CCM a livello Link Layer &
AES-256-GCM a livello applicativo \\

Nuovo materiale asimmetrico per ogni messaggio &
No &
No \\

Riconnessione con stato persistente &
Riutilizzo della LTK tramite bonding &
Riutilizzo della session key tramite resume \\

Overhead applicativo per messaggio &
Gestito principalmente dallo stack BLE &
37 B nel wire format PQ-BLE \\

Forward secrecy nelle riconnessioni &
Non garantita dal semplice riutilizzo della LTK &
Non garantita dal semplice riutilizzo della session key \\
Interazione utente & dipende dal pairing & SAS esplicito \\
Modifica stack BLE & no & no \\
Compatibilità & universale & solo dispositivi che implementano PQ-BLE \\
Stato del PoC & BLE standard maturo & validato come protocollo e
trasporto hardware \\
\end{longtable}

\subsection{Limiti progettuali}\label{limiti-progettuali}

\begin{enumerate}
\def\labelenumi{\arabic{enumi}.}
\tightlist
\item
  \textbf{Overhead radio maggiore}: ML-KEM usa chiavi e ciphertext molto
  più grandi di ECDH.
\item
  \textbf{Non è uno standard Bluetooth SIG}: è una soluzione applicativa
  custom.
\item
  \textbf{SAS non scala a grandi reti IoT}: richiede confronto umano.
\item
  \textbf{Firmware ancora demo transport}: il DK non esegue la
  crittografia completa.
\item
  \textbf{Nessuna cifratura link-layer}: SMP è disabilitato
  intenzionalmente; la protezione è applicativa.
\item
  \textbf{Side-channel non trattati}: un porting embedded reale dovrebbe
  considerare timing, cache, power analysis e RNG hardware.
\item
  \textbf{Notification raw nel firmware corrente}: la notification
  finale valida il trasporto, non la decrittazione AES-GCM on-device.
\end{enumerate}

\subsection{Perché non usare direttamente il Security Manager
BLE}\label{perchuxe9-non-usare-direttamente-il-security-manager-ble}

Integrare ML-KEM direttamente nello stack BLE richiederebbe modifiche al
Security Manager e quindi allo standard o a implementazioni specifiche
come Zephyr/BlueZ. Questo approccio sarebbe più pulito dal punto di
vista architetturale, ma molto più complesso e meno portabile.

PQ-BLE-HANDSHAKE sceglie invece un approccio applicativo:

\begin{lstlisting}
nessuna modifica BLE stack
compatibile con GATT standard
più facile da prototipare
meno integrato con link-layer security
\end{lstlisting}

Il trade-off è esplicito e coerente con un proof-of-concept
universitario.

\begin{center}\rule{0.5\linewidth}{0.5pt}\end{center}

\section{Conclusioni e sviluppi futuri}\label{capitolo-8---conclusioni-e-sviluppi-futuri}

\subsection{Conclusioni}\label{conclusioni}

Il progetto ha raggiunto l'obiettivo principale: dimostrare che è
possibile costruire un handshake post-quantum sopra BLE GATT senza
modificare lo stack BLE nativo.

I risultati principali sono:

\begin{itemize}
\tightlist
\item
  implementazione Python modulare del protocollo;
\item
  ML-KEM-768 integrato tramite liboqs;
\item
  SAS Numeric Comparison;
\item
  session key derivation tramite HKDF-SHA256;
\item
  canale AES-256-GCM con AAD e replay protection;
\item
  session resumption lato Python;
\item
  101 test automatizzati passanti;
\item
  firmware Zephyr per nRF54L15 DK compilato e flashato;
\item
  demo reale PC central $\leftrightarrow$ nRF54L15 DK completata;
\item
  notification raw da 57 byte ricevuta;
\item
  cattura Wireshark con nRF52840 Dongle che mostra il trasporto
  ATT/GATT.
\end{itemize}

La parte embedded completa, cioè ML-KEM/AES-GCM eseguiti direttamente
sul DK, rimane uno sviluppo futuro. Tuttavia, la validazione ottenuta è
sufficiente per dimostrare la fattibilità del trasporto BLE/GATT reale e
la correttezza del protocollo lato software.

Il principale overhead post-quantum è concentrato nel full handshake,
durante il quale vengono trasferiti la public key e il ciphertext
ML-KEM-768. Una volta derivata la session key, la comunicazione non
utilizza più primitive post-quantum costose, ma procede mediante
AES-256-GCM. Le riconnessioni possono evitare un nuovo scambio ML-KEM
attraverso la session resumption, al costo di una minore forward secrecy
dovuta alla persistenza della chiave.

\subsection{Sviluppi futuri}\label{sviluppi-futuri}

\begin{enumerate}
\def\labelenumi{\arabic{enumi}.}
\item
  \textbf{Porting ML-KEM su nRF54L15}\\
  Integrare una libreria PQC ottimizzata per Cortex-M33, misurando
  memoria, latenza e consumo.
\item
  \textbf{AES-GCM/HKDF on-chip}\\
  Derivare la session key e generare notification AES-GCM direttamente
  sul DK.
\item
  \textbf{Session store su flash}\\
  Salvare \passthrough{\lstinline!session\_id!} e session key in memoria
  persistente del microcontrollore.
\item
  \textbf{Autenticazione automatica con ML-DSA}\\
  Sostituire o affiancare il SAS con firme post-quantum.
\item
  \textbf{Handshake ibrido ECDH + ML-KEM}\\
  Combinare sicurezza classica e post-quantum in modo conservativo.
\item
  \textbf{Benchmark energetico}\\
  Misurare consumo radio e CPU su nRF54L15.
\item
  \textbf{Dissettore Wireshark e analisi automatica}\\
  Sviluppare un dissector o uno script di post-processing capace di riconoscere automaticamente header PQ-BLE, indici dei frammenti, comandi di controllo e notification, estendendo la cattura manuale già completata.
\item
  \textbf{Verifica formale del protocollo}\\
  Modellare PQ-BLE-HANDSHAKE con ProVerif o Tamarin nel modello simbolico
  di Dolev--Yao. Le proprietà da verificare includono segretezza della
  session key, agreement sulla trascrizione, autenticazione del peer dopo
  confronto SAS, resistenza al replay e separazione tra sessioni e ruoli.
\end{enumerate}

\begin{center}\rule{0.5\linewidth}{0.5pt}\end{center}

\appendix
\section{Estratti del codice critico}

Questa appendice riporta gli estratti minimi necessari a rendere la
relazione autosufficiente. Il codice completo è disponibile nella
repository pubblica del progetto.

\subsection{Derivazione del SAS}
\begin{lstlisting}[style=pythonstyle]
transcript = public_key + ciphertext + shared_secret
digest = hashlib.sha256(transcript).digest()
sas = int.from_bytes(digest[:4], "big") % 1_000_000
\end{lstlisting}

\subsection{AAD e wire format}
\begin{lstlisting}[style=pythonstyle]
aad = (
    session_id
    + sender_role
    + seq_num.to_bytes(8, "big")
    + msg_type
)
wire = seq_num_bytes + msg_type + iv + ciphertext_and_tag
\end{lstlisting}

\subsection{Resume limitato per tempo e utilizzi}
\begin{lstlisting}[style=pythonstyle]
session_key = store.load(session_id, increment_usage=False)
response = await wait_resume_response()
if response == RESUME_OK_NOTIFY:
    store.mark_used(session_id)
\end{lstlisting}

\subsection{Notification raw nella demo hardware}
\begin{lstlisting}[style=pythonstyle]
await channel.start_receiving(decrypt_notifications=False)
await client.write_control(b"START")
raw_payload = await channel.receive(timeout=5.0)
\end{lstlisting}

\section{Riproducibilità}

\subsection{Test automatici}
\begin{lstlisting}[style=bashstyle]
python -m pytest tests/ -v
# Risultato atteso: 101 passed
\end{lstlisting}

\subsection{Demo hardware}
\begin{lstlisting}[style=bashstyle]
python -m src.central.main --device PQ-BLE-Device --demo  --no-sas-confirm  --log-level DEBUG
\end{lstlisting}

\subsection{Build e flash del firmware}
\begin{lstlisting}[style=bashstyle]
west build -b nrf54l15dk/nrf54l15/cpuapp -p always
west flash
\end{lstlisting}

\subsection{Benchmark e aggiornamento automatico della tesina}
\begin{lstlisting}[style=bashstyle]
.\\benchmarks\\run_all.ps1
python scripts/generate_benchmark_latex.py
latexmk -pdf -interaction=nonstopmode \\
  report/tesina_pq_ble_handshake_finale_IT.tex
\end{lstlisting}

% \section{Evidenze e artefatti}

% Gli artefatti principali sono organizzati come segue:

% \begin{lstlisting}
% docs/hardware-validation-log.txt
% docs/images/*.png
% docs/captures/*.pcapng
% benchmarks/results/*.txt
% benchmarks/results/*.json
% report/tesina_pq_ble_handshake_finale_IT.tex
% report/tesina_pq_ble_handshake_finale_IT.pdf
% \end{lstlisting}

\section*{Bibliografia}
\addcontentsline{toc}{section}{Bibliografia}

\begin{enumerate}
\item NIST, \textbf{FIPS 203 - Module-Lattice-Based Key-Encapsulation
Mechanism Standard}, 2024. DOI: \url{https://doi.org/10.6028/NIST.FIPS.203}.
\item NIST, \textbf{FIPS 204 - Module-Lattice-Based Digital Signature
Standard}, 2024. DOI: \url{https://doi.org/10.6028/NIST.FIPS.204}.
\item NIST, \textbf{SP 800-227 - Recommendations for Key-Encapsulation
Mechanisms}, 2025.
\item NIST, \textbf{IR 8547 - Transition to Post-Quantum Cryptography
Standards}, Initial Public Draft, 2024.
\item Bluetooth SIG, \textbf{Bluetooth Core Specification 6.3}, Vol. 3,
Parts F, G e H, 2026.
\item H. Krawczyk e P. Eronen, \textbf{RFC 5869 - HMAC-based
Extract-and-Expand Key Derivation Function (HKDF)}, IETF, 2010.
\item NIST, \textbf{SP 800-38D - Recommendation for Block Cipher Modes of
Operation: Galois/Counter Mode (GCM) and GMAC}, 2007.
\item Open Quantum Safe, \textbf{liboqs: Open-Source Library for
Quantum-Safe Cryptographic Algorithms}, documentazione ufficiale.
\item H. Antila et al., \textbf{Bleak Documentation - Cross-platform
Asynchronous BLE GATT Client for Python}, documentazione ufficiale.
\item Zephyr Project, \textbf{Bluetooth Generic Attribute Profile (GATT)
APIs}, documentazione ufficiale.
\item Nordic Semiconductor, \textbf{nRF Sniffer for Bluetooth LE},
documentazione dello strumento di sviluppo.
\item Wireshark Foundation, \textbf{Wireshark User's Guide and Display
Filters Reference}.
\item B. Blanchet et al., \textbf{ProVerif: Automatic Cryptographic
Protocol Verifier - User Manual and Tutorial}, 2023.
\item S. Meier et al., \textbf{The Tamarin Prover for the Symbolic
Analysis of Security Protocols}, CAV, 2013; manuale aggiornato del Tamarin
Prover.
\item J. Bos et al., \textbf{CRYSTALS - Kyber: a CCA-Secure
Module-Lattice-Based KEM}, IEEE EuroS\&P, 2018.
\item O. Regev, \textbf{On Lattices, Learning with Errors, Random Linear
Codes, and Cryptography}, Journal of the ACM, 2009.
\end{enumerate}

\end{document}
